\section{A balanced closing statement can hide a stock shortage}
\label{sec:typed-economic-history}

The preceding history records monetary balances and movements. A company also
needs goods to deliver. A closing inventory value can agree with a closing
quantity even when the proposed sequence sells goods before they are available.
Testing that sequence requires quantities with their own units and conditions
on the intermediate states.

Consider an authored example with opening cash of 20 dollars, no inventory and
a credit equity balance of 20 dollars. Four units arrive at a fixed cost of
2 dollars each, and three units are sold for 3 dollars each. Both events are
dated 30 June 2024; the supplied order says which occurs first. The dates and
amounts are synthetic conformance inputs. The assumed postings illustrate a
fixed-cost, cash-settled transaction, without establishing the recognition
requirements of an actual inventory or revenue policy.

Let \(C\) denote cash in dollars, \(V\) inventory carrying amount in dollars,
\(E\) the signed, debit-positive equity balance in dollars, and \(Q\) the
number of units. Thus the opening state is
\((C,V,E,Q)=(20,0,-20,0)\). For a restock quantity \(u\) and a sale
quantity \(v\), the assumed effects are
\begin{align}
\label{eq:typed-stock-effects}
 (\Delta C,\Delta V,\Delta E,\Delta Q)_{\mathrm{restock}}
   &=(-2u,\;2u,\;0,\;u),\\
 (\Delta C,\Delta V,\Delta E,\Delta Q)_{\mathrm{sale}}
   &=(3v,\;-2v,\;-v,\;-v).
\end{align}
The coefficients 2 and 3 carry dollars per unit in the monetary components;
the quantity coefficients are dimensionless. The three monetary changes sum
to zero for each event. The quantity change is retained separately and is
never added to dollars merely because both are represented by rational numbers.

Substituting \(u=4\) first gives \((12,8,-20,4)\). Substituting \(v=3\)
next gives \((21,2,-23,1)\). The closing book balances because
\(21+2-23=0\), and the remaining unit has the assumed carrying amount 2.
Now reverse the order. The sale first gives \((29,-6,-23,-3)\), and the
later restock still produces \((21,2,-23,1)\). Every monetary posting
balances in both orders. Only the second order requires an intermediate stock
of minus three units.

\subsection{Apply a stated economic condition at every event}

Suppose the authored model permits between zero and 100 units at every state.
This is an explicit stock condition, not something inferred from the balance
sheet. Write \(x_j\) for the state just after event \(j\), with \(x_0\)
the opening, and \(Q_j\) for its quantity component. The conditional family
then requires
\begin{equation}
\label{eq:typed-stock-prefix}
 0\leq Q_j\leq100\qquad\text{for }j=0,1,\ldots,n.
\end{equation}
The restock-first path satisfies equation~\eqref{eq:typed-stock-prefix},
with quantities 0, 4 and 1. The sale-first path fails at its first event,
where \(Q_1=-3\). Testing only the final date would miss the violation.
Same-day events therefore retain their declared sequence, even though a daily
snapshot includes all events on that day.

The \path{TypedAffineHistory/1} reference permits a typed linear form of the
accounts and memory to be bounded in this way. Each supplied condition applies
to the opening and every post-event state. Its interval and proposed premise
are retained with the family. If no stock condition is supplied, the algebraic
sale-first path remains a candidate and the result reports that the condition
was not checked. A missing premise instead leaves the conditional question
unresolved. Neither case establishes that selling before delivery is allowed
for the company being studied.

The type of an event size need not be money. In equation~\eqref{eq:typed-stock-effects}
it is a quantity of goods, while each affected field has its own type. Before
forming scalar equations, the checker requires
\begin{equation}
\label{eq:typed-history-product}
 \operatorname{type}(b_{ij})\,\times\,
 \operatorname{type}(u_j)=\operatorname{type}(x_i),
\end{equation}
where \(b_{ij}\) is the fixed coefficient taking event size \(u_j\) into
field \(i\). Adding the coefficient's dimension exponents to those of the
event type, and multiplying their scales, must yield the field's exact dimension
and scale; their book
contexts must agree. For a cash effect, dollars per unit times units gives
dollars. A euro coefficient, a different book or an uncompensated thousand-unit
scale fails this check before solving. A measurement coefficient is checked
against its field and result type in the same way.

The reference retains the original types and definition after lowering to
scalar rational equations. Independent replay starts with the original opening
domains, applies the typed coefficients, checks every balanced posting and
declared intermediate condition, then recomputes the snapshot measurements.
Unmeasured economic memory remains a shared unknown. This extension is affine:
prices and other coefficients are fixed authored values, while event sizes and
openings may be unknown. Unknown prices multiplied by unknown quantities,
integer lots, endogenous event discovery and general contract methods need
additional methods. The declared stock condition checks one economic
restriction; it does not supply all operational or accounting prerequisites.

\subsection{Connect cash and quantity observations to the same path}
\label{sec:typed-history-observation-attachment}

The cash and stock observations can now restrict the same complete typed
history. This matters when neither restock nor sale quantity is supplied as a
point value. Retain the opening cash of 20 dollars and zero units, the fixed
cost of 2 dollars per unit, the sale price of 3 dollars per unit and the
restock-first sequence. Let both event quantities range from zero to ten units.
Suppose the observations report closing cash of 21 dollars and one remaining
unit. Opening cash, cash movement and closing cash form one temporal relation;
opening units, unit movement and closing units form another.

Write \(u\) for restocked units and \(v\) for sold units. The original
history's effects and the two closing observations jointly imply
\begin{align}
\label{eq:typed-observation-solution}
 u-v&=1,\\
 -2u+3v&=21-20=1.
\end{align}
Substituting \(u=v+1\) into the second equation gives
\(-2(v+1)+3v=1\), hence \(v=3\) and \(u=4\). Cash alone would leave
many possible pairs; the quantity observation supplies a distinct restriction.
Both equations apply to the same unknown events. Solving each parent separately
and passing an arbitrary witness from one into the other could discard their
common solution.

The closed \path{TypedHistoryMeasurementAttachment/1} record retains both
complete parent definitions and explicit mappings for the selected observations.
For an instant observation \(y_{k,t}\), the mapping has the form
\begin{equation}
\label{eq:typed-observation-mapping}
 y_{k,t}=\sum_i a_{ki}x_{i,t},\qquad
 \operatorname{type}(a_{ki})\times\operatorname{type}(x_i)
   =\operatorname{type}(y_k).
\end{equation}
Here \(x_{i,t}\) is a monetary account or economic-memory field at the
observation date and \(a_{ki}\) is a fixed authored coefficient. A unit
observation can use a dimensionless coefficient of one on unit stock. A dollar
valuation of that stock needs a stated dollars-per-unit coefficient. Every
product must yield the observation's exact dimensions, scale and book context;
the program supplies no implicit exchange rate, unit conversion or price.
The mapping's named quantity and proposed interpretation premise remain part
of its definition. Correct dimensions alone do not establish that a disclosure
and a stock record refer to the same population of goods.

For a period movement, the mapping instead sums the event changes in the exact
half-open period. It excludes the opening stock. If the example starts with
two units and still restocks four and sells three, the unit movement remains
one while closing stock is three. An instant snapshot includes every event
on that date in its declared order. It cannot select the state after restocking
but before a later same-day sale merely by using their shared calendar date.

The attachment includes the typed history's entire constraint system, not only
its scalar monetary lowering. In particular, every opening and post-event
stock bound remains. Reverse the example's event order: the same closing cash
and stock observations still agree, but the intermediate quantity is minus
three. The joined family rejects that path under the stated nonnegative-stock
condition. Omitting the condition leaves the algebraic candidate and reports
that the stock restriction was not checked.

Every observation and temporal relation remains present even when no mapping
is supplied. Such an observation creates an unresolved target obligation;
it does not disappear from its parent equations. Unknown economic memory also
remains free within its declared domains. Missing mapping or parent premises
prevent solving, while conflicting premise identities reject the definition.
A feasible assignment replays the original typed history and the original
temporal observation checker, recomputes the mappings with exact rational
arithmetic and satisfies every compiled constraint. The source interpretations
and complete actual company state still require their own evidence.

The retained examples include compatible alternative histories, individually
feasible parents whose intersection is empty, explicit valuation and scale
conversion, wrong-unit mappings and the intermediate shortage. Their complete
definitions and replay evidence are in
\path{data/accounting-master/typed-history-attachment-20261024}.
This connection constrains authored paths with observations. It does not yet
apply arbitrary accounting methods through changing legal and operational state.


\subsection{Author the same family as text and as typed records}

A conditional family must keep the same meaning when it is written as text.
The closed \path{ConditionalFamilyExchange/1} authoring interface covers joint
measurements, temporal measurements, monetary histories, typed affine histories
and the existing history--measurement attachment. Declarations name observations,
equations, accounts, memory, templates, events, measurements and intermediate
conditions. Their payloads retain the complete typed records, intervals,
premises and source references. The author supplies the computational budget
and the exact digest of the selected complete definition.

For example, the two event declarations in the stock case begin as follows;
the angle-bracketed domain markers here stand for the complete authored interval
records and are explanatory placeholders, not accepted syntax:
\begin{quote}\small\ttfamily
event \{"id":"restock","date":"2024-06-30",\\
\hspace*{1em}"template":"restock","amount":<four-unit domain>\};\\
event \{"id":"sale","date":"2024-06-30",\\
\hspace*{1em}"template":"sale","amount":<three-unit domain>\};
\end{quote}
Each complete declaration occupies one line in the actual source. The complete
examples contain the bounds, endpoint inclusion flags, named proposed premise
and provenance that the excerpt abbreviates. Reversing the two declarations
changes the event sequence and the selected definition's digest. It can change
feasibility under the stock condition even though the closing statement agrees.

Attachment declarations refer to earlier complete family definitions. The
parser embeds both parent formulas; it never substitutes their chosen witnesses.
The existing attachment still accepts the monetary-history and temporal-measurement
parents described in Section~\ref{sec:conditional-history-families}.
The separate typed attachment in Section~\ref{sec:typed-history-observation-attachment}
is not yet a profile of this five-profile exchange. Its complete typed records
therefore use the separate checked interface rather than that text grammar.

The text and byte decoders reject duplicate keys or declarations, unavailable
profiles, malformed UTF-8, floating or non-finite numeric tokens, wrong selection
digests, omitted budgets and trailing content. Rational quantities remain exact
reduced strings. Explicit zero budgets retain the original unknown result;
missing proposed premises remain unresolved. The parser treats declarations
as data and does not execute their contents or retrieve source references.

Map keys are normalized for both typed and text exchange inputs, while every
ordered list remains intact. This distinction matters because the preceding
solver reports assign internal variable names in insertion order: semantically
equal maps can otherwise produce different diagnostic text. Replay uses the
normalized definition, and independent witness checks also agree with the
original unsorted record. This normalization does not promise identical outcomes
for arbitrary old requests stopped at tight resource limits.

The full grammar, strict error behavior and eight retained source-to-record-to-result
examples accompany the implementation. The monetary embedding also preserves
the earlier alternative histories, and adverse cases retain the stock shortage,
missing premises and resource stops. The supporting specifications and checks
are under \path{data/accounting-master/}:
\begin{quote}\small
\path{typed-affine-history-20261024}\\
\path{conditional-family-authoring-20261024}
\end{quote}
These references make the declared state and authoring subset reproducible.
Whole-company method coverage, richer contracts, source admission and an
independent implementation of the broader language remain separate tasks.
